\documentclass[11pt,a4paper]{article}
\usepackage[T1]{fontenc}
\usepackage[utf8]{inputenc}
\usepackage{amsmath,amssymb}
\usepackage{booktabs}
\usepackage{siunitx}
\usepackage{microtype}
\usepackage{geometry}
\usepackage{hyperref}
\usepackage{cleveref}
\usepackage{multirow}
\usepackage{array}

\geometry{margin=25mm}
\sisetup{per-mode=symbol, inter-unit-product=\ensuremath{\cdot}}

\newcommand{\DG}{\ensuremath{D_\gamma}}
\newcommand{\DA}{\ensuremath{D_\alpha}}
\newcommand{\Xrex}{\ensuremath{X_\mathrm{rex}}}
\newcommand{\epsr}{\ensuremath{\varepsilon_\mathrm{ret}}}
\newcommand{\epsrate}{\ensuremath{\dot\varepsilon}}
\newcommand{\rhod}{\ensuremath{\rho_\mathrm{d}}}
\newcommand{\sy}{\ensuremath{\sigma_\mathrm{y}}}
\newcommand{\fv}{\ensuremath{f_\mathrm{v}}}
\newcommand{\Nv}{\ensuremath{N_\mathrm{v}}}
\newcommand{\Ae}{\ensuremath{\mathrm{Ae}_3}}

\title{MICDEL: A Coupled Pass-by-Pass Model for Microstructure \\
  Evolution During Hot Rolling of Microalloyed Steels}
\author{Implementation report --- \texttt{micdeltool}}
\date{March 2026}

\begin{document}
\maketitle

\begin{abstract}
We describe the mathematical formulation and Modula-3 implementation of a
MICDEL-type model for predicting microstructure evolution during multi-pass hot
rolling of microalloyed steels.  The model couples static recrystallization
(JMAK), normal grain growth with Zener drag, classical nucleation and growth of
MX precipitates (VN, VC, V(C,N), NbC, Nb(C,N), TiN, TiC), the
austenite-to-ferrite transformation, and yield stress prediction via
superposition of lattice friction, solid solution, Hall--Petch, Orowan--Ashby,
and Taylor dislocation hardening.  Comprehensive tests across eight steel
compositions---from plain C--Mn to multi-microalloyed V--Ti--N---demonstrate
physically consistent trends in grain refinement, strain accumulation, and
precipitation strengthening.  The model is based on the
Lagneborg--Hutchinson--Siwecki--Zajac monograph.
\end{abstract}

\tableofcontents

%=============================================================================
\section{Introduction}
%=============================================================================

The thermomechanical controlled processing (TMCP) of microalloyed steels
exploits the interplay between deformation, recrystallization, precipitation,
and phase transformation to achieve fine ferrite grain sizes and high yield
strengths without expensive alloying or post-rolling heat treatment.  Predicting
the outcome of a given rolling schedule requires a coupled, physically-based
model that tracks all relevant microstructural variables pass by pass.

The MICDEL (MICrostructure DEveLogement) approach, as described in the
monograph by Lagneborg, Hutchinson, Siwecki, and Zajac~\cite{lagneborg}, provides
such a framework.  This report describes our Modula-3 implementation of the core
MICDEL model, presents the governing equations, and validates the model against
eight steel compositions spanning the principal microalloying strategies.

%=============================================================================
\section{Model overview}
%=============================================================================

The state vector tracked through each rolling pass comprises:
\begin{itemize}
  \item Austenite grain diameter \DG{} (\si{\metre})
  \item Recrystallized fraction \Xrex{} (0--1)
  \item Retained strain \epsr{}
  \item Dislocation density \rhod{} (\si{\per\metre\squared})
  \item Temperature $T$ (\si{\kelvin})
  \item Dissolved solute concentrations (\si{wt.\%}): C, N, V, Nb, Ti, Mn, Si
  \item Precipitate populations: $\{\Nv, \bar{r}, \fv\}$ for each species
\end{itemize}

The simulation proceeds as follows:
\begin{enumerate}
  \item \textbf{Initialization}: Set \DG{} and composition from reheat conditions.
  \item \textbf{For each pass}:
    \begin{enumerate}
      \item Apply deformation (update \epsr, \rhod).
      \item Evolve interpass interval (coupled time-stepping of recrystallization,
            grain growth, and precipitation).
    \end{enumerate}
  \item \textbf{Transformation}: Cool to \Ae, compute ferrite grain size \DA,
        simulate coil cooling (ferrite precipitation), predict yield stress \sy.
\end{enumerate}

%=============================================================================
\section{Thermodynamics}
\label{sec:thermo}
%=============================================================================

\subsection{Solubility products}

The solubility of a precipitate species MX in a given matrix phase is expressed
as
\begin{equation}
  \log_{10}\bigl([\mathrm{M}]\cdot[\mathrm{X}]\bigr) \;=\; -\frac{A}{T} + B
  \label{eq:solprod}
\end{equation}
where $[\mathrm{M}]$ and $[\mathrm{X}]$ are the dissolved concentrations of the
metallic and interstitial elements in weight percent, and $T$ is the absolute
temperature in Kelvin.  The constants $A$ and $B$ are specific to each
precipitate species and matrix phase.

\begin{table}[htbp]
\centering
\caption{Solubility product parameters from the Lagneborg monograph.}
\label{tab:solprod}
\begin{tabular}{@{}llSS@{}}
\toprule
Species & Phase & {$A$ (K)} & {$B$} \\
\midrule
VN     & Austenite & 7840  & 3.02 \\
VC     & Austenite & 6560  & 4.45 \\
V(C,N) & Austenite & 7700  & 3.30 \\
NbC    & Austenite & 6770  & 2.26 \\
Nb(C,N)& Austenite & 7500  & 2.50 \\
TiN    & Austenite & 15790 & 5.40 \\
TiC    & Austenite & 7000  & 2.75 \\
\midrule
VN     & Ferrite   & 8330  & 3.90 \\
VC     & Ferrite   & 7050  & 4.29 \\
V(C,N) & Ferrite   & 8100  & 3.80 \\
NbC    & Ferrite   & 7500  & 3.00 \\
Nb(C,N)& Ferrite   & 8000  & 3.20 \\
TiN    & Ferrite   & 15790 & 5.40 \\
TiC    & Ferrite   & 7500  & 3.10 \\
\bottomrule
\end{tabular}
\end{table}

The equilibrium solubility product at temperature $T$ is:
\begin{equation}
  K_\mathrm{s}(T) = 10^{-A/T + B}
\end{equation}

\subsection{Driving force for precipitation}

The volume free energy driving nucleation and growth of species MX is
\begin{equation}
  \Delta G_\mathrm{v} = \frac{RT}{V_\mathrm{m}}
    \ln\!\left(\frac{c_\mathrm{M}\,c_\mathrm{X}}{K_\mathrm{s}(T)}\right)
  \label{eq:driving}
\end{equation}
where $R = \SI{8.314}{\joule\per\mole\per\kelvin}$ is the gas constant and
$V_\mathrm{m}$ is the molar volume of the precipitate.  When
$c_\mathrm{M}\,c_\mathrm{X} \le K_\mathrm{s}(T)$, the system is in solid
solution and $\Delta G_\mathrm{v} = 0$.

\subsection{Diffusivity}

The diffusivity of the rate-controlling solute follows an Arrhenius form:
\begin{equation}
  D(T) = D_0 \exp\!\left(-\frac{Q_\mathrm{d}}{RT}\right)
  \label{eq:diffusivity}
\end{equation}

\begin{table}[htbp]
\centering
\caption{Diffusion and interfacial energy data.}
\label{tab:diffusion}
\begin{tabular}{@{}lSSSSS@{}}
\toprule
Species & {$\gamma$ (\si{\joule\per\metre\squared})}
  & {$V_\mathrm{m}$ (\SI{e-6}{\metre\cubed\per\mole})}
  & {$D_0^{\gamma}$ (\SI{e-5}{\metre\squared\per\second})}
  & {$Q_\mathrm{d}^{\gamma}$ (\si{\kilo\joule\per\mole})} \\
\midrule
VN      & 0.5  & 10.29 & 3.7  & 186 \\
VC      & 0.5  & 10.86 & 3.7  & 186 \\
V(C,N)  & 0.5  & 10.50 & 3.7  & 186 \\
NbC     & 0.5  & 13.38 & 8.3  & 266 \\
Nb(C,N) & 0.5  & 13.00 & 8.3  & 266 \\
TiN     & 0.8  & 11.63 & 1.5  & 250 \\
TiC     & 0.8  & 12.18 & 1.5  & 250 \\
\bottomrule
\end{tabular}
\end{table}

%=============================================================================
\section{Precipitation kinetics}
\label{sec:precip}
%=============================================================================

Each precipitate species is tracked as a mean-radius population
$\{\Nv,\,\bar r,\,\fv\}$ representing the number density, mean radius, and
volume fraction respectively.

\subsection{Classical nucleation}

The steady-state nucleation rate is
\begin{equation}
  J = N_0\, Z\, \beta^* \exp\!\left(-\frac{\Delta G^*}{k_\mathrm{B} T}\right)
  \label{eq:nucleation}
\end{equation}
where the critical free energy barrier is
\begin{equation}
  \Delta G^* = \frac{16\pi\,\gamma^3}{3\,\Delta G_\mathrm{v}^2}
  \label{eq:dgstar}
\end{equation}
The pre-factors are:
\begin{align}
  N_0 &= 10^{28} + \rhod \times 10^{10}
    & &\text{(nucleation site density, m$^{-3}$)}
    \label{eq:n0} \\[6pt]
  Z &= \frac{V_\mathrm{m}}{2\pi\gamma}
       \sqrt{\frac{\Delta G_\mathrm{v}}{N_\mathrm{A}\,k_\mathrm{B} T}}
    & &\text{(Zeldovich factor)}
    \label{eq:zeldovich} \\[6pt]
  \beta^* &= \frac{4\pi\,r^{*2}\,D\,c_\mathrm{M}/100}
                   {\bigl(V_\mathrm{m}/N_\mathrm{A}\bigr)^2}
    & &\text{(atomic impingement rate)}
    \label{eq:betastar}
\end{align}
with the critical radius
\begin{equation}
  r^* = \frac{2\gamma}{\Delta G_\mathrm{v}}
  \label{eq:rstar}
\end{equation}

The heterogeneous nucleation on dislocations is captured through the
$\rhod$-dependent term in \cref{eq:n0}: higher dislocation densities
(from retained strain after incomplete recrystallization) provide more
nucleation sites and accelerate precipitation.

\subsection{Diffusion-controlled growth}

The growth rate of an existing precipitate of radius $r$ is
\begin{equation}
  \frac{\mathrm{d}r}{\mathrm{d}t} = \frac{D}{r}\,
    \frac{c_0 - c_\mathrm{i}}{c_\mathrm{p} - c_\mathrm{i}}
  \label{eq:growth}
\end{equation}
where $c_0$ is the far-field (matrix) concentration of the metallic solute,
$c_\mathrm{p} \approx 1$ is the concentration in the precipitate (taken as
pure MX), and $c_\mathrm{i}$ is the interfacial equilibrium concentration
modified by the Gibbs--Thomson capillarity correction:
\begin{equation}
  c_\mathrm{i} = \frac{K_\mathrm{s}(T)}{c_\mathrm{X}}\,
    \exp\!\left(\frac{2\gamma V_\mathrm{m}}{r\,RT}\right)
  \label{eq:gibbs-thomson}
\end{equation}

\subsection{Ostwald ripening (LSW)}

When precipitates exceed twice the critical radius, coarsening follows the
Lifshitz--Slyozov--Wagner (LSW) theory:
\begin{equation}
  \frac{\mathrm{d}(\bar r^3)}{\mathrm{d}t}
    = \frac{8\,\gamma\,D\,c_\mathrm{i}\,V_\mathrm{m}}{9\,RT}
  \label{eq:lsw}
\end{equation}
which gives an effective radial growth rate:
\begin{equation}
  \frac{\mathrm{d}\bar{r}}{\mathrm{d}t}\bigg|_\mathrm{coarsen}
    = \frac{1}{3\bar{r}^2}\,\frac{\mathrm{d}(\bar r^3)}{\mathrm{d}t}
\end{equation}

\subsection{Solute depletion and stoichiometric cap}

As precipitates form, dissolved solute is consumed.  The maximum attainable
volume fraction is limited by stoichiometry:
\begin{equation}
  f_{\mathrm{v},\max} \approx \min(c_\mathrm{M},\, c_\mathrm{X}) \times 0.02
  \label{eq:fvmax}
\end{equation}
where the factor 0.02 arises from the ratio of precipitate to matrix densities
and molar masses for typical MX carbides and nitrides.

The solute depletion from a volume fraction increment $\Delta\fv$ is
\begin{equation}
  \Delta c_\mathrm{M} = \frac{\Delta\fv}{0.02}
  \label{eq:depletion}
\end{equation}
For carbonitrides V(C,N) and Nb(C,N), the interstitial depletion is split
equally between C and N.

\subsection{Mean-radius update}

At each timestep, the mean radius is updated as a weighted average of the
existing population (which has grown) and newly nucleated precipitates (at the
critical radius):
\begin{equation}
  \bar{r}_\mathrm{new} = \frac{\Nv^\mathrm{old}\,(\bar{r} + \Delta r)
    + \Delta\Nv \cdot r^*}
    {\Nv^\mathrm{old} + \Delta\Nv}
  \label{eq:rmean}
\end{equation}
The volume fraction is then
\begin{equation}
  \fv = \tfrac{4}{3}\pi\,\bar{r}^3\,\Nv
  \label{eq:fv}
\end{equation}

%=============================================================================
\section{Static recrystallization}
\label{sec:rex}
%=============================================================================

Static recrystallization (SRX) during the interpass interval follows the JMAK
(Johnson--Mehl--Avrami--Kolmogorov) kinetics.

\subsection{Half-recrystallization time}

The time for 50\% recrystallization is given by:
\begin{equation}
  t_{0.5} = A\,\varepsilon^{-p}\,\epsrate^{-q}\,D_0^{\,s}\,
    \exp\!\left(\frac{Q_\mathrm{rex}}{RT}\right) \cdot f_\mathrm{drag}
  \label{eq:t50}
\end{equation}
where $D_0$ is in micrometers and the default
parameters are:
\begin{equation*}
  A = \num{2.5e-19},\quad p = 4.0,\quad q = 1.0,\quad s = 2.0,\quad
  Q_\mathrm{rex} = \SI{230}{\kilo\joule\per\mole}
\end{equation*}

\subsection{Solute drag retardation}

The presence of dissolved Nb and V retards recrystallization through solute drag
on grain boundaries:
\begin{equation}
  f_\mathrm{drag} = \exp\!\left(\frac{275\,000\,c_\mathrm{Nb}}{100}\right)
    \cdot \exp\!\left(\frac{3000\,c_\mathrm{V}}{100}\right)
  \label{eq:solutedrag}
\end{equation}
where $c_\mathrm{Nb}$ and $c_\mathrm{V}$ are in weight percent.  The much
larger coefficient for Nb reflects its far stronger solute drag effect
compared to V, which is a key reason why Nb steels exhibit strain
accumulation (pancaking) while V steels recrystallize freely between passes.

\subsection{JMAK recrystallized fraction}

The Avrami equation gives the recrystallized fraction as a function of time:
\begin{equation}
  \Xrex(t) = 1 - \exp\!\left[-0.693\left(\frac{t}{t_{0.5}}\right)^{\!n}\right]
  \label{eq:jmak}
\end{equation}
with Avrami exponent $n = 1.0$.

\subsection{Recrystallized grain size}

The grain size after complete recrystallization follows Sellars:
\begin{equation}
  D_\mathrm{rex}\;(\si{\micro\metre})
    = 0.743\,D_0^{\,0.67}\,\varepsilon^{-0.67}
  \label{eq:drex}
\end{equation}
where $D_0$ is the prior austenite grain size in micrometers.  During partial
recrystallization, the effective grain size is the mixture:
\begin{equation}
  \DG = \Xrex\,D_\mathrm{rex} + (1 - \Xrex)\,D_\gamma^\mathrm{old}
  \label{eq:dmix}
\end{equation}

%=============================================================================
\section{Grain growth and Zener drag}
\label{sec:gg}
%=============================================================================

Normal grain growth in the recrystallized austenite follows power-law kinetics
with a Zener drag limit imposed by second-phase particles.

\subsection{Grain growth kinetics}

The grain growth rate from parabolic ($m = 2$) kinetics is:
\begin{equation}
  \frac{\mathrm{d}D}{\mathrm{d}t}
    = \frac{k_0\,\exp(-Q_\mathrm{gg}/RT)}{m\,D^{m-1}}
  \label{eq:gg}
\end{equation}
with $k_0 = \SI{4.1e-3}{\metre\squared\per\second}$ and
$Q_\mathrm{gg} = \SI{200}{\kilo\joule\per\mole}$.

\subsection{Zener limiting grain size}

The maximum grain size permitted by precipitate pinning is:
\begin{equation}
  D_\mathrm{lim} = \frac{\alpha\,\bar{r}_\mathrm{Z}}{\fv^\mathrm{total}}
  \label{eq:zener}
\end{equation}
where $\alpha = 0.05$ is the Zener constant, $\bar{r}_\mathrm{Z}$ is the
volume-fraction-weighted mean precipitate radius across all species, and
$\fv^\mathrm{total}$ is the total precipitate volume fraction.

This Zener limit is the principal mechanism by which TiN particles---which
are extremely stable and present from reheating---control the austenite
grain size.

%=============================================================================
\section{Deformation model}
\label{sec:deform}
%=============================================================================

Upon each roll pass with strain $\varepsilon$ and strain rate $\epsrate$:
\begin{align}
  \epsr &\leftarrow \epsr + \varepsilon
    \label{eq:epsacc} \\
  \rhod &\leftarrow \rhod + \varepsilon\,\epsrate \times 10^{13}
    \label{eq:rhoincr}
\end{align}
with an upper cap of $\rhod \le 10^{16}\;\si{\per\metre\squared}$.  The
recrystallized fraction is reset to zero:
\begin{equation}
  \Xrex \leftarrow 0
\end{equation}

As recrystallization proceeds during the interpass interval, the dislocation
density and retained strain are reduced proportionally to the recrystallized
fraction increment:
\begin{align}
  \rhod &\leftarrow \rhod\,(1 - 0.9\,\Delta\Xrex) \\
  \epsr &\leftarrow \epsr\,(1 - \Delta\Xrex)
\end{align}

%=============================================================================
\section{Coupled interpass evolution}
\label{sec:interpass}
%=============================================================================

The core of the model is the time-stepping loop during each interpass interval.
With timestep $\Delta t$ (default \SI{0.01}{\second}):

\begin{enumerate}
  \item \textbf{Cool}: $T \leftarrow T - \dot{T}_\mathrm{cool}\,\Delta t$
  \item \textbf{Recrystallize}: Compute $t_{0.5}$ from current state; update
        \Xrex{} via JMAK; mix grain sizes; reduce \rhod{} and \epsr.
  \item \textbf{Grain growth}: If $\Xrex > 0.5$, grow \DG{} subject to Zener
        limit from current precipitate population.
  \item \textbf{Precipitate}: For each species (VN, VC, V(C,N), NbC, Nb(C,N),
        TiN, TiC), compute nucleation, growth, and coarsening; deplete solute.
\end{enumerate}

The feedback between precipitation and recrystallization is critical:
\begin{itemize}
  \item Precipitates exert Zener drag, limiting grain growth.
  \item Dissolved solute (especially Nb) retards recrystallization through
        solute drag.
  \item Retained strain from incomplete recrystallization provides dislocation
        nucleation sites, accelerating strain-induced precipitation.
\end{itemize}

%=============================================================================
\section{Austenite-to-ferrite transformation}
\label{sec:transform}
%=============================================================================

\subsection{\texorpdfstring{$\Ae$}{Ae3} temperature}

The equilibrium transformation temperature follows the Andrews formula:
\begin{equation}
  \Ae\;(\mathrm{K}) = 1456.15 - 203\sqrt{c_\mathrm{C}}
    - 30\,c_\mathrm{Mn} + 44.7\,c_\mathrm{Si}
    - 15.2\,c_\mathrm{Ni} - 11.0\,c_\mathrm{Cr}
  \label{eq:ae3}
\end{equation}

\subsection{Ferrite grain size (Sellars--Beynon)}

The ferrite grain size is predicted by the Sellars--Beynon model:
\begin{equation}
  \DA\;(\si{\micro\metre}) = 3.75 + 0.18\,\DG\;(\si{\micro\metre})
    + \frac{0.29}{\sqrt{\dot{T}_\mathrm{cool}}}
    - 0.7\,c_\mathrm{Mn}
  \label{eq:dalpha}
\end{equation}
with an additional refinement factor from intragranular nucleation on V(C,N)
particles:
\begin{equation}
  \DA \leftarrow \frac{\DA}{1 + 10^{-23}\,\Nv^{\mathrm{VN}}\,\bar{r}^{\mathrm{VN}}
    + 10^{-23}\,\Nv^{\mathrm{VCN}}\,\bar{r}^{\mathrm{VCN}}}
  \label{eq:igr}
\end{equation}

\subsection{Coil cooling}

After transformation, the model simulates isothermal coil cooling at
\SI{0.01}{\kelvin\per\second} for up to \SI{3600}{\second}, during which
V(C,N) precipitation occurs in the ferrite phase.  Carbon is partitioned
to pearlite (ferrite solubility capped at \SI{0.02}{wt.\%}).

%=============================================================================
\section{Yield stress prediction}
\label{sec:strength}
%=============================================================================

The yield stress is a linear superposition of five strengthening contributions:
\begin{equation}
  \boxed{\;
  \sy = \sigma_0 + \sigma_\mathrm{ss} + \sigma_\mathrm{HP}
    + \sigma_\mathrm{Or} + \sigma_\mathrm{d}
  \;}
  \label{eq:sigma}
\end{equation}

\subsection{Lattice friction}

\begin{equation}
  \sigma_0 = \SI{53.9}{\mega\pascal}
\end{equation}

\subsection{Solid solution strengthening}

\begin{equation}
  \sigma_\mathrm{ss}\;(\mathrm{MPa})
    = 4570\,c_\mathrm{C} + 3750\,c_\mathrm{N}
    + 32\,c_\mathrm{Mn} + 84\,c_\mathrm{Si} + 11\,c_\mathrm{Mo}
  \label{eq:ss}
\end{equation}
where all concentrations are in weight percent.

\subsection{Hall--Petch grain refinement}

\begin{equation}
  \sigma_\mathrm{HP} = \frac{k_\mathrm{y}}{\sqrt{D_\alpha}}
    \quad\text{with}\quad k_\mathrm{y}
    = 18.14\;\mathrm{MPa\!\cdot\!mm}^{1/2}
  \label{eq:hp}
\end{equation}
where $D_\alpha$ is in mm.

\subsection{Orowan--Ashby precipitation hardening}

For each precipitate species with volume fraction $\fv$ and mean radius
$\bar{r}$, the interparticle spacing is:
\begin{equation}
  \lambda = \left(\sqrt{\frac{2\pi}{3\fv}}
    - \sqrt{\frac{8}{3}}\right)\bar{r}
  \label{eq:lambda}
\end{equation}
and the Orowan stress increment is:
\begin{equation}
  \sigma_\mathrm{Or} = \sum_i
    \frac{0.538\,G\,b}{\lambda_i}\,\ln\!\left(\frac{\bar{r}_i}{b}\right)
  \label{eq:orowan}
\end{equation}
where $G = \SI{81000}{\mega\pascal}$ is the ferrite shear modulus and
$b = \SI{2.48e-10}{\metre}$ is the Burgers vector.  The sum runs over all
precipitate species with non-negligible volume fractions.

The Orowan--Ashby mechanism is the primary strengthening route exploited in
V--N steels: dense, fine V(C,N) precipitates formed during coil cooling
provide a large $\sigma_\mathrm{Or}$ contribution.

\subsection{Taylor dislocation hardening}

\begin{equation}
  \sigma_\mathrm{d} = \alpha_\mathrm{T}\,M\,G\,b\,\sqrt{\rhod}
  \label{eq:taylor}
\end{equation}
with Taylor factor $M = 3.06$ and $\alpha_\mathrm{T} = 0.33$.  After
transformation, $\rhod$ is reset to $10^{12}\;\si{\per\metre\squared}$,
giving a negligible contribution ($< 1\;\mathrm{MPa}$) for typical steels.

%=============================================================================
\section{Implementation}
\label{sec:impl}
%=============================================================================

The model is implemented in Modula-3 \cite{m3} as two packages within the
\texttt{m3utils} tree:

\begin{description}
  \item[\texttt{micdellib}] --- Library containing nine modules:
    \begin{itemize}
      \item \texttt{MicdelThermo}: Solubility products, diffusion data,
            driving force (\cref{sec:thermo})
      \item \texttt{MicdelAlloy}: Steel composition record and accessors
      \item \texttt{MicdelPrecipitate}: Mean-radius population kinetics
            (\cref{sec:precip})
      \item \texttt{MicdelRecryst}: JMAK recrystallization (\cref{sec:rex})
      \item \texttt{MicdelGrainGrowth}: Grain growth with Zener drag
            (\cref{sec:gg})
      \item \texttt{MicdelTransform}: $\gamma\to\alpha$ transformation
            (\cref{sec:transform})
      \item \texttt{MicdelStrength}: Yield stress prediction
            (\cref{sec:strength})
      \item \texttt{MicdelState}: Coupled state vector and interpass evolution
            (\cref{sec:interpass})
      \item \texttt{MicdelSchedule}: Schedule file parser
    \end{itemize}
  \item[\texttt{micdeltool}] --- Driver program that reads a schedule file,
    runs the pass-by-pass simulation, and outputs results.
\end{description}

\subsection{Units and dimensional conventions}

All internal calculations use SI units (meters, Kelvin,
seconds, Pascals).  Many empirical correlations in the metallurgical
literature are calibrated in engineering units (micrometers for grain size,
millimeters for Hall--Petch).  Unit conversions are applied at the boundaries
of each empirical formula:
\begin{itemize}
  \item \textbf{Recrystallization} (\cref{eq:t50,eq:drex}):
    $D_0\,(\mathrm{m}) \to D_0\,(\mathrm{\mu m})$ via
    $D_0^{\mu\mathrm{m}} = D_0 \times 10^6$ on entry;
    $D_\mathrm{rex}\,(\mathrm{\mu m}) \to D_\mathrm{rex}\,(\mathrm{m})$
    via $\times 10^{-6}$ on exit.
  \item \textbf{Ferrite grain size} (\cref{eq:dalpha}):
    same $\mathrm{m} \leftrightarrow \mathrm{\mu m}$ conversion.
  \item \textbf{Hall--Petch} (\cref{eq:hp}):
    $D_\alpha\,(\mathrm{m}) \to D_\alpha\,(\mathrm{mm})$ via $\times 10^3$.
  \item \textbf{Solute concentrations}: uniformly in wt.\%\ throughout.
\end{itemize}

\subsection{Explicit Euler time integration}
\label{sec:euler}

The interpass evolution (\cref{sec:interpass}) is integrated using the
explicit (forward) Euler method.  Given the state vector
$\mathbf{u}^{(n)}$ at time $t_n$, the update is:
\begin{equation}
  \mathbf{u}^{(n+1)} = \mathbf{u}^{(n)}
    + \Delta t\;\mathbf{f}\!\left(\mathbf{u}^{(n)},\, t_n\right)
  \label{eq:euler}
\end{equation}
where $\mathbf{f}$ encapsulates the coupled right-hand sides (temperature
cooling, recrystallization kinetics, grain growth, and precipitation for all
seven species).  The default timestep is $\Delta t = \SI{0.01}{\second}$.

Concretely, at each Euler step the following sub-updates are applied
sequentially:

\paragraph{Step 1: Temperature.}
\begin{equation}
  T^{(n+1)} = T^{(n)} - \dot{T}_\mathrm{cool}\,\Delta t
\end{equation}

\paragraph{Step 2: Recrystallization.}
Compute $t_{0.5}$ from the current $\epsr$, $\DG$, $T$, and solute state
using \cref{eq:t50}.  Evaluate the JMAK fraction (\cref{eq:jmak}) at elapsed
time $t_\mathrm{elapsed} + \Delta t$.  If $\Xrex$ increased by
$\Delta \Xrex > 0$:
\begin{align}
  D_\mathrm{rex} &= 0.743\,D_\gamma^{\,0.67}\,\varepsilon_\mathrm{ret}^{-0.67}
    \quad\text{(\cref{eq:drex}, in $\mu$m)} \\
  \DG^{(n+1)} &= \Xrex^{(n+1)}\,D_\mathrm{rex}
    + (1 - \Xrex^{(n+1)})\,\DG^{(n)} \\
  \rhod^{(n+1)} &= \rhod^{(n)}\,(1 - 0.9\,\Delta\Xrex) \\
  \epsr^{(n+1)} &= \epsr^{(n)}\,(1 - \Delta\Xrex)
\end{align}
The factor 0.9 in the dislocation reduction reflects that recrystallization
does not fully annihilate all dislocations; a residual density of
$10^{10}\;\mathrm{m}^{-2}$ is enforced as a floor.

\paragraph{Step 3: Grain growth.}
Applied only when $\Xrex > 0.5$ (i.e., the material is predominantly
recrystallized).  The grain size increment from \cref{eq:gg} is:
\begin{equation}
  \Delta D = \frac{k_0\,e^{-Q_\mathrm{gg}/RT}}{m\,D^{m-1}}\,\Delta t
\end{equation}
subject to the Zener ceiling $D \le D_\mathrm{lim}$ from \cref{eq:zener}.
The update is monotonic: grain size never decreases from grain growth alone.

\paragraph{Step 4: Precipitation (for each of the 7 species).}
For species $s$ with current population $\{\Nv,\,\bar{r},\,\fv\}$:
\begin{enumerate}
  \item Check if solute is available and $\fv < f_{\mathrm{v},\max}$.
  \item Compute the nucleation rate $J$ from \cref{eq:nucleation}:
    $\Delta\Nv = J\,\Delta t$.
  \item Compute the growth rate $\mathrm{d}r/\mathrm{d}t$ from \cref{eq:growth}
    and the coarsening rate from \cref{eq:lsw}:
    $\Delta r = (\mathrm{d}r/\mathrm{d}t|_\mathrm{grow}
    + \mathrm{d}r/\mathrm{d}t|_\mathrm{coarsen})\,\Delta t$.
  \item Update the mean radius via the weighted average (\cref{eq:rmean}).
  \item Compute the new volume fraction $\fv = \frac{4}{3}\pi\bar{r}^3\Nv$.
  \item If $\fv > f_{\mathrm{v},\max}$, cap $\fv$ and adjust $\Nv$ to be
    consistent: $\Nv = \fv / (\frac{4}{3}\pi\bar{r}^3)$.
  \item Deplete solute: $c_\mathrm{M} \leftarrow c_\mathrm{M}
    - \Delta\fv / 0.02$.
\end{enumerate}

\noindent This operator-splitting approach (temperature $\to$ recrystallization
$\to$ grain growth $\to$ precipitation) introduces a first-order splitting
error of $O(\Delta t)$, consistent with the Euler integration order.  The
default $\Delta t = 0.01\;\mathrm{s}$ is small relative to typical interpass
intervals (3--10\,s), giving $\sim$300--1000 steps per interpass and adequate
resolution of the kinetics.

\subsection{Transformation and coil cooling}
\label{sec:impl-transform}

The post-rolling transformation is handled in three stages:

\begin{enumerate}
  \item \textbf{Cool to $\Ae$}: If the current temperature exceeds the
    equilibrium $\Ae$ (\cref{eq:ae3}), advance time by
    $\Delta t_\mathrm{cool} = (T - \Ae)/\dot{T}_\mathrm{cool}$ and set
    $T = \Ae$.

  \item \textbf{Compute ferrite grain size}: Apply the Sellars--Beynon
    formula (\cref{eq:dalpha}) with the Mn correction and intragranular
    refinement factor (\cref{eq:igr}).  Partition carbon to pearlite by
    capping ferrite $c_\mathrm{C}$ at 0.02\,wt.\%.  Reset $\rhod$ to
    $10^{12}\;\mathrm{m}^{-2}$ (transformation-annealed state).

  \item \textbf{Coil cooling loop}: Iterate for up to 3600 steps of
    $\Delta t = 1\;\mathrm{s}$, cooling at the coil rate
    (\SI{0.01}{\kelvin\per\second}).  At each step, run the precipitation
    sub-update (Step~4 above) for all species in the \emph{ferrite} phase,
    using ferrite solubility products and diffusion data.  Exit if
    $T < \SI{300}{\degreeCelsius}$.
\end{enumerate}

\noindent Finally, the yield stress is computed from \cref{eq:sigma} using the
post-coil-cooling state.

\subsection{Nucleation rate capping}

The classical nucleation rate (\cref{eq:nucleation}) involves the exponential
$\exp(-\Delta G^*/k_\mathrm{B}T)$, which can produce extremely large or
extremely small values depending on supersaturation.  The implementation
applies a hard cap of $J \le 10^{30}\;\mathrm{m}^{-3}\mathrm{s}^{-1}$ to
prevent numerical overflow, and floors $J$ at zero (no negative nucleation).

\subsection{Input file format}

Schedules are specified in a simple text format:
\begin{verbatim}
COMPOSITION  C=0.09  N=0.013  V=0.09  Ti=0.01  Mn=1.4  Si=0.3
REHEAT       T=1523  D_gamma=50e-6
PASS  eps=0.3  rate=10  T=1373  t_inter=5  cool=2
PASS  eps=0.3  rate=10  T=1343  t_inter=5  cool=2
...
COOLING  rate=15  T_coil=873  t_coil=3600  coilRate=0.01
\end{verbatim}
The parser reads key\texttt{=}value pairs from each line.  Temperatures are
in Kelvin, grain sizes in meters, times in seconds, and cooling rates in
K/s.  Lines beginning with \texttt{\#} are comments.  The parser collects
\texttt{PASS} lines into a dynamically-sized array using a linked list
during parsing, then converts to a contiguous array for the simulation loop.

\subsection{Software architecture}
\label{sec:architecture}

The complete implementation comprises approximately 1620 lines of Modula-3
source code, organized as a library (\texttt{micdellib}, 9~modules, 1367~LOC)
and a driver program (\texttt{micdeltool}, 1~module, 146~LOC), plus
88~lines across eight test schedule files.  \Cref{tab:loc} shows the
breakdown by module.

\begin{table}[htbp]
\centering
\caption{Lines of code by module (interfaces + implementations).}
\label{tab:loc}
\begin{tabular}{@{}llr@{}}
\toprule
Module & Role & LOC \\
\midrule
\texttt{MicdelThermo}       & Solubility products, diffusion data        & 145 \\
\texttt{MicdelAlloy}        & Composition record, element access          &  66 \\
\texttt{MicdelPrecipitate}  & Nucleation, growth, coarsening, depletion  & 296 \\
\texttt{MicdelRecryst}      & JMAK kinetics, solute drag, rex'd grain size & 96 \\
\texttt{MicdelGrainGrowth}  & Parabolic growth, Zener limit              &  81 \\
\texttt{MicdelTransform}    & Ae3, Sellars--Beynon, coil cooling         & 128 \\
\texttt{MicdelStrength}     & 5-component yield stress                   &  89 \\
\texttt{MicdelState}        & State vector, interpass coupling, transform & 198 \\
\texttt{MicdelSchedule}     & Text file parser                           & 268 \\
\midrule
\texttt{Main}               & Driver program                             & 146 \\
\midrule
\multicolumn{2}{@{}l}{Total (library + driver)} & 1513 \\
\bottomrule
\end{tabular}
\end{table}

The module dependency graph reflects the physical layering of the model.
At the bottom sit the ``leaf'' modules with no intra-library dependencies:

\begin{itemize}
  \item \texttt{MicdelThermo} --- pure data (solubility products, diffusion
    coefficients, interfacial energies) and stateless query functions.
    No mutable state.
  \item \texttt{MicdelAlloy} --- a record type holding the steel composition
    in wt.\%, with element accessors.
\end{itemize}

\noindent The kinetics modules form the middle tier, each importing
\texttt{Thermo} and \texttt{Alloy}:

\begin{itemize}
  \item \texttt{MicdelPrecipitate} imports \texttt{Thermo} and \texttt{Alloy}.
    Its \texttt{Step} procedure advances one species' population by one
    timestep, depleting solute from the \texttt{Alloy.T} record passed by
    \texttt{VAR} reference.  This is the largest module (296~LOC) because it
    handles all three kinetic regimes (nucleation, growth, coarsening) plus
    the stoichiometric cap and solute depletion bookkeeping.
  \item \texttt{MicdelRecryst} imports \texttt{Thermo} and \texttt{Alloy}.
    Stateless: computes $t_{0.5}$, $X$, and $D_\mathrm{rex}$ from the current
    state without side effects.
  \item \texttt{MicdelGrainGrowth} imports \texttt{Precipitate} (for the
    Zener limit, which requires the full precipitate population array).
  \item \texttt{MicdelTransform} imports \texttt{Alloy} and
    \texttt{Precipitate} (for intragranular refinement from VN particles).
  \item \texttt{MicdelStrength} imports \texttt{Alloy} and
    \texttt{Precipitate}.
\end{itemize}

\noindent The coupling module \texttt{MicdelState} sits at the top of the
library.  It imports all five kinetics modules plus \texttt{Thermo} and
\texttt{Alloy}, and provides the three entry points that implement the
simulation loop:

\begin{enumerate}
  \item \texttt{ApplyDeformation} --- accumulates strain and updates
    dislocation density after a roll pass (no time integration).
  \item \texttt{EvolveInterpass} --- the Euler time-stepping loop described
    in \cref{sec:euler}, calling \texttt{Recryst}, \texttt{GrainGrowth}, and
    \texttt{Precipitate.Step} (for each of 7~species) at every timestep.
  \item \texttt{Transform} --- invokes \texttt{MicdelTransform} for the
    $\gamma\to\alpha$ conversion and \texttt{MicdelStrength} for the final
    yield stress, plus a coil cooling sub-loop that calls
    \texttt{Precipitate.Step} in the ferrite phase.
\end{enumerate}

\noindent Finally, \texttt{MicdelSchedule} is an I/O module that parses the
text schedule format into structured records; it imports \texttt{Alloy} (for
the composition record) but is otherwise independent of the physics modules.

The driver program (\texttt{Main.m3}) imports only \texttt{MicdelState} and
\texttt{MicdelSchedule}.  Its logic is a straightforward loop:
\begin{verbatim}
  state := MicdelState.Init(sched.comp, sched.D_gamma_0, sched.T_reheat);
  FOR i := 0 TO LAST(sched.passes^) DO
    MicdelState.ApplyDeformation(state, pass.strain, pass.strainRate);
    MicdelState.EvolveInterpass(state, pass.t_interpass, pass.coolingRate);
    print one-line summary of state
  END;
  MicdelState.Transform(state, sched.cooling);
  print final results and yield stress breakdown
\end{verbatim}

\noindent External dependencies are minimal: the standard library
\texttt{libm3} (for \texttt{Math}, \texttt{Fmt}, \texttt{Wr}, \texttt{Rd},
\texttt{Params}, \texttt{Scan}), plus \texttt{cit\_util} for
\texttt{LRVector.T} (a dynamically-sized \texttt{LONGREAL} array used in
the schedule parser).  No numerical solvers, linear algebra packages, or
optimization libraries are required---the explicit Euler scheme needs only
elementary arithmetic.

The complete source code is publicly available at:\\
\url{https://github.com/mikanystrom/intel-async/tree/claude_edits/async-toolkit/m3utils/m3utils/micdel}

\subsection{Timing benchmarks}
\label{sec:timing}

\Cref{tab:timing} reports wall-clock execution times for each test schedule,
measured as the average of 100 consecutive runs on an Apple~M4~Pro (ARM64\_DARWIN).
Process startup overhead (Modula-3 runtime initialization, schedule file
parsing) accounts for approximately \SI{7}{\milli\second}; the actual
simulation computation ranges from \SI{1.5}{\milli\second} for the simpler
schedules to \SI{4}{\milli\second} for the multi-species cases.

\begin{table}[htbp]
\centering
\caption{Wall-clock execution times (average of 100 runs, Apple~M4~Pro).}
\label{tab:timing}
\begin{tabular}{@{}lSS@{}}
\toprule
Steel & {Total (ms)} & {Computation (ms)} \\
\midrule
Plain C--Mn       &  8.3 & 1.6 \\
V--N              &  8.7 & 2.0 \\
Nb                &  9.5 & 2.8 \\
Nb--V             & 10.4 & 3.7 \\
Ti-only           &  8.3 & 1.6 \\
V--Ti--N (RCR)    &  8.7 & 2.0 \\
V--Ti--N (CR)     &  9.8 & 3.1 \\
High-V            & 10.5 & 3.8 \\
\bottomrule
\end{tabular}
\end{table}

\noindent The sub-\SI{5}{\milli\second} computation times make the model
suitable for rapid parameter studies, schedule optimization loops, or
integration into online process control systems where thousands of schedule
evaluations may be required.

%=============================================================================
\section{Test results}
\label{sec:results}
%=============================================================================

Eight steel compositions were tested, spanning the principal microalloying
strategies employed in industrial TMCP practice.

\subsection{Steel compositions}

\begin{table}[htbp]
\centering
\caption{Test steel compositions (wt.\%).  All steels use a 6-pass schedule
  with reheat at \SI{1250}{\degreeCelsius} (\SI{1523}{\kelvin}).}
\label{tab:comp}
\begin{tabular}{@{}lccccccc@{}}
\toprule
Steel & C & N & V & Nb & Ti & Mn & Si \\
\midrule
Plain C--Mn       & 0.12  & 0.005 & ---  & ---  & ---   & 1.20 & 0.25 \\
V--N              & 0.08  & 0.015 & 0.12 & ---  & ---   & 1.30 & 0.30 \\
Nb                & 0.07  & 0.005 & ---  & 0.04 & ---   & 1.40 & 0.25 \\
Nb--V             & 0.08  & 0.012 & 0.08 & 0.03 & ---   & 1.40 & 0.30 \\
Ti-only           & 0.10  & 0.008 & ---  & ---  & 0.015 & 1.30 & 0.25 \\
V--Ti--N (RCR)    & 0.09  & 0.013 & 0.09 & ---  & 0.01  & 1.40 & 0.30 \\
V--Ti--N (CR)     & 0.09  & 0.013 & 0.09 & ---  & 0.01  & 1.40 & 0.30 \\
High-V            & 0.06  & 0.020 & 0.15 & ---  & 0.01  & 1.50 & 0.30 \\
\bottomrule
\end{tabular}
\end{table}

The RCR schedules finish above
$T_\mathrm{nr}$ (no-recrystallization temperature) with entry temperatures
from \SI{1100}{\degreeCelsius} down to \SI{950}{\degreeCelsius},
while the CR schedule uses the same V--Ti--N composition but finishes in the
non-recrystallization regime at \SI{750}{\degreeCelsius}.

\subsection{Summary of final results}

\begin{table}[htbp]
\centering
\caption{Predicted microstructure and mechanical properties.}
\label{tab:results}
\begin{tabular}{@{}lSSSSSSS@{}}
\toprule
Steel & {\DG} & {\Xrex} & {\DA}
  & {$\sigma_\mathrm{ss}$} & {$\sigma_\mathrm{HP}$}
  & {$\sigma_\mathrm{Or}$} & {\sy} \\
  & {(\si{\micro\metre})} & {} & {(\si{\micro\metre})}
  & {(\si{\mega\pascal})} & {(\si{\mega\pascal})}
  & {(\si{\mega\pascal})} & {(\si{\mega\pascal})} \\
\midrule
Plain C--Mn    & 17.3  & 1.000 & 6.1   & 169.6 & 231.9 &   0.0 & 455 \\
V--N           & 15.5  & 1.000 & 5.7   & 158.2 & 240.3 &   4.0 & 456 \\
Nb             & 60.0  & 0.000 & 13.6  & 157.2 & 155.3 &  32.3 & 399 \\
Nb--V          & 50.0  & 0.000 & 11.8  & 161.4 & 166.7 &  75.5 & 458 \\
Ti-only        &  9.8  & 1.000 &  4.7  & 154.0 & 264.7 &  79.6 & 552 \\
V--Ti--N (RCR) & 14.9  & 1.000 &  5.5  & 161.4 & 243.9 &  87.3 & 547 \\
V--Ti--N (CR)  & 10.8  & 0.869 &  4.8  & 161.4 & 262.2 & 442.6 & 920 \\
High-V         & 15.3  & 1.000 &  5.5  & 164.6 & 243.9 &  64.2 & 527 \\
\bottomrule
\end{tabular}
\end{table}

\noindent All steels include a lattice friction stress
$\sigma_0 = 53.9\;\mathrm{MPa}$ and negligible dislocation
hardening ($\sigma_\mathrm{d} < 1\;\mathrm{MPa}$) after transformation.

\subsection{Precipitate populations}

\begin{table}[htbp]
\centering
\caption{Final precipitate populations after coil cooling.}
\label{tab:precip}
\begin{tabular}{@{}llSSS@{}}
\toprule
Steel & Species
  & {$\Nv$ (\si{\per\metre\cubed})}
  & {$\bar{r}$ (\si{\nano\metre})}
  & {$\fv$} \\
\midrule
Plain C--Mn & --- & {---} & {---} & {---} \\
\midrule
V--N        & V(C,N)  & 2.4e15  & 473   & 1.07e-3 \\
\midrule
Nb          & Nb(C,N) & 1.0e19  & 21    & 4.0e-4 \\
\midrule
\multirow{2}{*}{Nb--V}
            & V(C,N)  & 2.1e17  & 121   & 1.6e-3 \\
            & Nb(C,N) & 2.1e22  & 0.7   & 2.9e-5 \\
\midrule
Ti-only     & TiN     & 1.1e24  & 0.3   & 1.6e-4 \\
\midrule
\multirow{2}{*}{V--Ti--N (RCR)}
            & V(C,N)  & 2.6e14  & 942   & 9.0e-4 \\
            & TiN     & 1.5e24  & 0.3   & 2.0e-4 \\
\midrule
\multirow{2}{*}{V--Ti--N (CR)}
            & V(C,N)  & 1.2e24  & 0.6   & 9.9e-4 \\
            & TiN     & 1.5e24  & 0.3   & 2.0e-4 \\
\midrule
\multirow{2}{*}{High-V}
            & V(C,N)  & 6.5e13  & 1436  & 8.0e-4 \\
            & TiN     & 1.9e24  & 0.3   & 2.0e-4 \\
\bottomrule
\end{tabular}
\end{table}

\subsection{Pass-by-pass evolution: V--Ti--N RCR}

\Cref{tab:passdata} shows the full pass-by-pass evolution for the V--Ti--N
RCR steel, demonstrating how the microstructural state develops
through the rolling sequence.

\begin{table}[htbp]
\centering
\caption{Pass-by-pass state for V--Ti--N steel (RCR schedule).}
\label{tab:passdata}
\begin{tabular}{@{}cSSSSS@{}}
\toprule
Pass & {$T$ (\si{\degreeCelsius})} & {\DG{} (\si{\micro\metre})}
  & {\Xrex} & {\epsr} & {$\fv^\mathrm{total}$} \\
\midrule
1 & 1090 & 22.9 & 1.000 & 0.000 & 2.0e-4 \\
2 & 1060 & 13.6 & 1.000 & 0.000 & 2.0e-4 \\
3 & 1032 & 14.8 & 1.000 & 0.000 & 5.4e-4 \\
4 & 1002 & 15.6 & 1.000 & 0.000 & 9.1e-4 \\
5 &  974 & 15.6 & 1.000 & 0.000 & 1.1e-3 \\
6 &  944 & 14.9 & 1.000 & 0.000 & 1.1e-3 \\
\bottomrule
\end{tabular}
\end{table}

Complete recrystallization occurs between all passes ($\Xrex = 1.0$), as
expected for V steels where solute drag is weak.  The austenite grain size
stabilizes around \SI{15}{\micro\metre} due to the balance between
recrystallization refinement and Zener-limited grain growth from TiN particles.
The total precipitate volume fraction increases from passes 3--5 as V(C,N)
begins to form in the austenite.

%=============================================================================
\section{Discussion}
\label{sec:discussion}
%=============================================================================

\subsection{Plain C--Mn baseline}

With no microalloy additions, the plain C--Mn steel shows complete
recrystallization between all passes ($\Xrex = 1.0$) and no precipitate
pinning.  The final austenite grain size of \SI{17.3}{\micro\metre}
progressively refines through successive recrystallization cycles, giving a
ferrite grain size of \SI{6.1}{\micro\metre} and yield stress of
\SI{455}{\mega\pascal} dominated by solid solution (\SI{170}{\mega\pascal})
and Hall--Petch (\SI{232}{\mega\pascal}) contributions.

\subsection{V--N steel: precipitation strengthening}

The V--N steel also recrystallizes fully (V has minimal solute drag), giving a
similar grain size to the baseline.  The V(C,N) precipitates that form during
austenite rolling have coarsened to a mean radius of \SI{473}{\nano\metre} by
the time of transformation---too large for significant Orowan strengthening.
Only \SI{4}{\mega\pascal} of Orowan hardening results.  The modest yield stress
of \SI{456}{\mega\pascal} reflects that the strengthening from V(C,N)
precipitation relies on very fine particles formed during coil cooling, which in
this simulation have time to coarsen during the \SI{3600}{\second} coil hold.

\subsection{Nb steel: strain accumulation}

The Nb steel demonstrates the classic pancaking (controlled rolling) behavior.
The strong solute drag from dissolved Nb
($f_\mathrm{drag} = \exp(275\,000 \times c_\mathrm{Nb}/100)$ in \cref{eq:solutedrag})
completely suppresses recrystallization below
$T_\mathrm{nr}$: $\Xrex = 0$ throughout all six passes.  The retained strain
accumulates to $\epsr = 1.5$, and the austenite grain remains at its initial
\SI{60}{\micro\metre}.  The lower yield
stress of \SI{399}{\mega\pascal} compared to the baseline reflects the larger
effective ferrite grain size (\SI{13.6}{\micro\metre}).  The Nb(C,N) precipitates
contribute only \SI{32}{\mega\pascal} of Orowan hardening.

In industrial practice, the pancaked austenite microstructure provides extensive
grain boundary area for ferrite nucleation, which our Sellars--Beynon model
(\cref{eq:dalpha}) captures only approximately; a more detailed treatment
would account for the effective boundary area per unit volume of the
deformed austenite, producing a finer ferrite grain size.

\subsection{Nb--V combined microalloying}

Combining Nb and V provides both strain accumulation (from Nb solute drag) and
precipitation strengthening (from V(C,N)).  The Nb suppresses recrystallization
as before, while V(C,N) precipitates form during coil cooling with a mean radius
of \SI{121}{\nano\metre} and $\fv = 1.6 \times 10^{-3}$.  The Orowan contribution of
\SI{76}{\mega\pascal} partially compensates for the larger ferrite grain size,
yielding \SI{458}{\mega\pascal}---comparable to the plain C--Mn steel but with
lower carbon and better toughness.

\subsection{Ti-only: grain refinement through Zener drag}

The Ti-only steel exploits the exceptional stability of TiN precipitates
($A = 15\,790$, the largest among all species).  TiN forms at very high
temperatures and remains essentially undissolved throughout rolling, providing
persistent Zener drag that limits the austenite grain size to
\SI{9.8}{\micro\metre}---the smallest among the RCR steels.  This translates
to a fine ferrite grain of \SI{4.7}{\micro\metre} and a Hall--Petch contribution
of \SI{265}{\mega\pascal}, giving the highest RCR yield stress of
\SI{552}{\mega\pascal}.

Note that the TiN precipitates themselves also contribute \SI{80}{\mega\pascal}
of Orowan hardening, though in practice TiN particles are typically too coarse
for significant precipitation strengthening; this reflects the mean-field
approximation assigning all particles the same mean radius.

\subsection{V--Ti--N: RCR vs.\ CR comparison}

The comparison between RCR and CR schedules for the same V--Ti--N composition
is particularly instructive.

In the \textbf{RCR} schedule (all passes above $T_\mathrm{nr}$), complete
recrystallization occurs between passes.  V(C,N) precipitates that form in
austenite have time to coarsen to $\bar{r} = \SI{942}{\nano\metre}$, giving a
modest Orowan contribution of \SI{87}{\mega\pascal}.  Total yield stress:
\SI{547}{\mega\pascal}.

In the \textbf{CR} schedule (last four passes at 900--\SI{750}{\degreeCelsius}),
a dramatically different precipitate population emerges.  The lower temperatures
produce an enormous nucleation density of V(C,N)
($\Nv = 1.2 \times 10^{24}\;\si{\per\metre\cubed}$) with very fine mean radius
(\SI{0.6}{\nano\metre}).  The resulting Orowan contribution of
\SI{443}{\mega\pascal} dominates the total yield stress of
\SI{920}{\mega\pascal}.

This extreme value highlights the theoretical potential of controlled rolling
with V--N steels, though the predicted Orowan increment likely overstates what
would be achieved in practice: the sub-nanometer precipitate sizes approach the
classical-to-atomistic transition where the mean-field nucleation and growth
model breaks down.  Nevertheless, the qualitative trend---that CR produces much
finer V(C,N) dispersions and stronger precipitation hardening than RCR---is
well established experimentally.

\subsection{High-V steel}

Increasing V to \SI{0.15}{wt.\%} with \SI{0.020}{wt.\%} N provides more solute
for precipitation.  However, under RCR conditions the V(C,N) precipitates
coarsen to $\bar{r} = \SI{1436}{\nano\metre}$---the largest observed---giving
only \SI{64}{\mega\pascal} of Orowan hardening.  More V means more total
precipitate, but also more solute available for growth and coarsening.  The
lesson is clear: high V content alone does not guarantee high strengthening; the
rolling strategy and cooling conditions must be matched to the alloy design.

\subsection{General observations}

The results demonstrate several well-known principles of microalloyed steel
design:
\begin{enumerate}
  \item \textbf{Nb for recrystallization control}: The exponential solute drag
    from Nb is the most effective mechanism for suppressing SRX and accumulating
    strain.
  \item \textbf{Ti for grain refinement}: TiN's extreme thermodynamic stability
    provides robust Zener pinning throughout the entire hot rolling temperature
    range.
  \item \textbf{V for precipitation strengthening}: V(C,N) precipitation in
    ferrite during coil cooling is the primary strengthening mechanism, but its
    effectiveness depends critically on the precipitate size distribution.
  \item \textbf{The rolling strategy matters}: The same V--Ti--N steel gives
    \SI{547}{\mega\pascal} (RCR) versus \SI{920}{\mega\pascal} (CR), a
    \SI{370}{\mega\pascal} difference entirely from processing.
\end{enumerate}

%=============================================================================
\section{Conclusion}
%=============================================================================

A coupled, physically-based MICDEL model has been implemented in Modula-3,
comprising nine library modules and a driver program.  The model successfully
captures the essential physics of microalloyed steel hot rolling:
solubility-limited precipitation, classical nucleation theory,
diffusion-controlled growth and LSW coarsening, JMAK static recrystallization
with solute drag, Zener-limited grain growth, Sellars--Beynon ferrite grain size
prediction, and superposition yield stress modeling.

Testing across eight steel compositions demonstrates physically consistent
trends and realistic yield stress predictions in the
400--\SI{550}{\mega\pascal} range for conventional RCR schedules, with
the CR V--Ti--N case illustrating the extreme potential of fine precipitation
strengthening.

\begin{thebibliography}{9}
\bibitem{lagneborg}
  R.~Lagneborg, T.~Siwecki, S.~Zajac, and B.~Hutchinson,
  ``The Role of Vanadium in Microalloyed Steels,''
  \emph{Scandinavian Journal of Metallurgy}, vol.~28, pp.~186--241, 1999.

\bibitem{sellars}
  C.~M. Sellars and J.~A. Whiteman,
  ``Recrystallization and grain growth in hot rolling,''
  \emph{Metal Science}, vol.~13, pp.~187--194, 1979.

\bibitem{ashby}
  M.~F. Ashby,
  ``The deformation of plastically non-homogeneous materials,''
  \emph{Philosophical Magazine}, vol.~21, pp.~399--424, 1970.

\bibitem{gladman}
  T.~Gladman,
  ``Precipitation hardening in metals,''
  \emph{Materials Science and Technology}, vol.~15, pp.~30--36, 1999.

\bibitem{zajac}
  S.~Zajac, T.~Siwecki, W.~B. Hutchinson, and R.~Lagneborg,
  ``Strengthening mechanisms in vanadium microalloyed steels intended for
  long products,''
  \emph{ISIJ International}, vol.~38, pp.~1130--1139, 1998.

\bibitem{andrews}
  K.~W. Andrews,
  ``Empirical formulae for the calculation of some transformation
  temperatures,''
  \emph{Journal of the Iron and Steel Institute}, vol.~203, pp.~721--727, 1965.

\bibitem{sellars-beynon}
  C.~M. Sellars and J.~H. Beynon,
  ``Microstructural development during hot rolling,''
  in \emph{Proc.\ Conf.\ on High Strength Low Alloy Steels},
  Wollongong, Australia, pp.~142--150, 1984.

\bibitem{m3}
  L.~Cardelli, J.~Donahue, L.~Glassman, M.~Jordan, B.~Kalsow, and G.~Nelson,
  ``Modula-3 language definition,''
  \emph{ACM SIGPLAN Notices}, vol.~27, no.~8, pp.~15--42, 1992.
\end{thebibliography}

\end{document}
